\def\udk{УДК 004.9}





\def\tit{МЕТОД РЕГИСТРАЦИИ МОРФИНГА ТРЕХМЕРНОГО ОБЪЕКТА НА~ОСНОВЕ 


ДАННЫХ НАТУРНОГО ЭКСПЕРИМЕНТА}





  \def\aut{О.\,П.~Архипов$^1$, Ю.\,А.~Маньяков$^2$, Д.\,О.~Сиротинин$^3$}





\titel{\udk}{\tit}{\aut}





\def\razd{Метод регистрации морфинга трехмерного объекта\ldots} 


%данных натурного эксперимента}





\def\sect{О.\,П.~Архипов, Ю.\,А.~Маньяков, Д.\,О.~Сиротинин}





\index{Архипов О.\,П.}


\index{Маньяков Ю.\,А.}


\index{Сиротинин Д.\,О.}





\renewcommand{\thefootnote}{\arabic{footnote}}


\footnotetext[1]{ОФ ИПИ РАН, ofran@orel.ru}


\footnotetext[2]{ОФ ИПИ РАН, maniakov\_yuri@mail.ru }


\footnotetext[3]{ОФ ИПИ РАН, vespert@mail.ru}








\Abst{Описан метод регистрации изменения формы натурных  объектов, 


основанный на анализе и обработке двумерных изображений, полученных оптическими 


регистриру\-ющи\-ми приборами с использованием методов фотограмметрии.\linebreak Метод 


позволяет регистрировать изменения натурного объекта, строить трехмерную модель 


данного объекта и отображать на ней изменения натурного объекта.}





\KW{стереоизображения; цветовые характеристики; трехмерный объект; регистрация; 


морфинг} 





 \thispagestyle{empty}





\section{Введение}





На сегодняшний день актуальна проблема получения и представления информации о 


морфинге (изменении формы) трех\-мер\-ных объектов. Системы, позволяющие получать 


подобную\linebreak информа\-цию, могут быть использованы для наблюдения за динамикой 


различных физических объектов и анализа их поведения,  в системах мониторинга на 


производстве, геомониторинга и в медицине, для удаленного управления 


ро\-бо\-та\-ми-ма\-ни\-пу\-ля\-то\-ра\-ми в различных агрессивных средах, а также для 


обучения их операторов.





Основными методами получения информации о поведении\linebreak трехмерных объектов, 


использующимися на сегодняшний день являются:


\begin{itemize}


\item видеомониторинг;


\item различные системы регистрации кинематики (в частности, технология motion 


capture);


\item аналитическое описание морфинга.


\end{itemize}





Видеомониторинг~[1] основан на использовании видео- или реже~--- фотооборудования. 


Данный метод имеет ряд недостатков, а именно: отсутствие объективного контроля, 


недостаточно точен, не позволяет получать данные об объекте в числовом виде 


(координаты), требует наличие оператора, не позволяет представлять информацию в 


трехмерном виде, затруднено предо\-став\-ле\-ние информации в реальном времени в связи с 


не\-об\-хо\-ди\-мостью передавать большие объемы данных.





Существующие системы регистрации кинематики можно\linebreak условно разделить на две 


группы: маркерные (регистрация кинематики ведется с помощью отслеживания 


перемещений особых маркерных точек, наносимых тем или иным образом на 


ис-\linebreak следуемый объект) и безмаркерные (основаны на фото- или\linebreak видеорегистрирующей 


методике, где маркерные точки и кинематические \mbox{звенья} отмечаются вручную либо 


автоматически  непосредственно на фото- или видеокадрах). Маркерные системы 


являются лабораторными, т.\,е.\ для их реализации необходимы определенные условия.





Существует довольно много различных типов маркерных\linebreak систем (оптические, магнитные, 


механические, гироскопиче-\linebreak ские), однако все они обладают схожими недостатками: 


\begin{itemize}


\item системы подразумевают размещение на исследуемом объекте различных 


регистрирующих датчиков, подчас до\-воль\-но громоздких и затрудняющих перемещения, а 


также не имеющие возможности размещения на определенные объекты;


\item подвержены помехам, зашумлению, порождающим неточности измерений;


\item все системы довольно дорогостоящи;


\item результирующая информация, полученная в результате работы этих систем, требует 


дополнительной обработки.


\end{itemize}





Безмаркерные системы подразумевают необходимость видеосъемки с разных ракурсов 


для получения точной ха\-рак\-те\-ри\-сти\-ки движения, и имеют невысокую точность 


измерений.





Существуют достаточно много аналитических моделей описывающих изменение формы и 


поведение различных объектов, однако, все они прежде всего основываются на 


математических моделях, различных алгоритмах прогнозирования и имеют ряд 


недостатков, среди которых большая вычислительная слож-\linebreak ность, и недостаточная 


гибкость.





С учетом недостатков описанных систем с \textit{целью} ускорения и упрощения процесса 


регистрации морфинга трехмерных объектов предлагается метод, позволяющий 


отказаться от физического нанесения маркеров в пользу идентификации виртуальных 


маркеров~--- опорных точек (ОТ) на плоских изображениях (фотографиях, кадрах 


видеосъемки), цветовые характеристики, которых служат для идентификации 


составляющих их пикселей с последующей регистрацией перемещений последних в 


пространстве. 





\section{Описание метода регистрации морфинга на~основе данных натурного 


эксперимента}





Метод можно подразделить на следующие этапы: 


\begin{enumerate}[1.]


\item Регистрация поведения объекта при помощи фото- либо видеоаппаратуры.


\item Аппроксимация цветового пространства изображений.


\item Сегментация изображений по цветовым характеристикам\linebreak пикселей.


\item Структурирование сегментированных областей.


\item Фиксация опорных точек на плоских изображениях.


\item Вычисление координат ОТ в трехмерном пространстве по двумерным координатам 


на плоских изображениях.


\item Построение и преобразование модели (морфинг) в соответствии с изменениями 


координат ОТ.


\end{enumerate}





На рис.~1 отображена общая схема технологии, где 


Cl~--- цветовые характеристики 


изображения; CIPL, CIPR~--- цветовая индексная палитра для левого и правого кадров; 


CSL, CSR~--- сегментированные области для левого и правого кадров; СOPL,\linebreak СOPR~--- 


ограничивающие прямоугольники для левого и правого кадров; COT~--- координаты 


ОТ в двумерном пространстве; V~--- воксельная модель; $B_1\ldots B_m$~--- сетка 


ограничи-\linebreak вающих объемов; $D_1\ldots D_k$~--- векторы смещения ОТ; $N$~---\linebreak мат\-ри\-ца 


смежности; $\{\mathrm{OT}_1\ldots 


\mathrm{OT}_n\}_1\ldots\{\mathrm{OT}_1\ldots\mathrm{OT}_n\}_k$~--- координаты ОТ в 


трехмерном пространстве для всех шагов\linebreak морфинга; ОО~--- ограничивающие объемы; 


$C_1\ldots C_n$~--- цветовые характеристики опорных точек.


      


\begin{figure} %fig1-7


%\vspace*{6pt}


\begin{center}


\mbox{%


\epsfxsize=109.78mm


\epsfbox{arh-1.eps}


}


 \end{center}


 \vspace*{-12pt}


\Caption{Технологическая схема метода регистрации мор\-фин\-га трехмерного объекта 


на основе данных натурного эксперимента 


\label{f1-7}}


\end{figure}





\subsection{Регистрация поведения объекта при~помощи фото- либо видеоаппаратуры} 


%2.1





Рассмотрим работу метода регистрации поведения объекта при использовании 


фотоаппаратуры. Восстановление трехмерных координат по двумерным изображениям 


осуществляется на основании бинокулярного восприятия объекта~[2]. В~процессе съемки 


создается пара изображений, левое и правое, получаемые с камер расположенных на 


расстоянии базисной величины друг от друга. Такой вид съемки называется 


стереосъемкой. Оси объектива камеры направлены параллельно друг другу. Выбор 


значения стереобазиса основывается на размерах объекта и расстоянии до него~[3]. 


Натурный объект располагается на экспериментальном стенде, имеющем специальные 


маркеры, один\linebreak из которых выбирается центром начала локальной системы координат, 


используемой для расчета трехмерных координат. Получение изображений всей 


поверхности объекта может осуществляться либо последовательным перемещение пары 


камер\linebreak вокруг объекта или установкой дополнительных камер, либо путем установки 


стенда с объектом на специальной поворачивающейся подставке~[4, 5].





Так как целью наблюдения за объектом является реконструкция его поведения во 


времени, то процесс съемки происходит через определенные промежутки времени. 


Данный временной\linebreak показатель ограничивается аппаратными возможностями оп-\linebreak тического 


регистрирующего прибора или частотой изменений\linebreak самого объекта. В~связи с этим 


исходные изображения формируются на каждом временном этапе. Формально каждое 


изображение является совокупностью пикселей, характеризующихся цветовыми 


координатами:


      \begin{gather*}


      A_1=


      \begin{bmatrix}


      R_{1,1} & R_{1,2} &\ldots & R_{1,m}\\


      R_{2,1} & R_{2,2} &\ldots &R_{2,m}\\


      \ldots &\ldots &\ldots &\ldots\\


      R_{n,1} & R_{n,2} & \ldots & R_{n,m}


      \end{bmatrix}\,;\quad 


      A_2=


      \begin{bmatrix}


      G_{1,1} & G_{1,2} &\ldots & G_{1,m}\\


      G_{2,1} & G_{2,2} &\ldots & G_{2,m}\\


      \ldots &\ldots &\ldots &\ldots\\


      G_{n,1} & G_{n,2} & \ldots & G_{n,m}


      \end{bmatrix}\,;\\


      A_3=


      \begin{bmatrix}


      B_{1,1} & B_{1,2} &\ldots & B_{1,m}\\


      B_{2,1} & B_{2,2} &\ldots & B_{2,m}\\


      \ldots &\ldots &\ldots &\ldots\\


      B_{n,1} & B_{n,2} & \ldots & B _{n,m}


      \end{bmatrix}\,,


      \end{gather*}


где  $R,G,B \in \{0,1,2,\ldots ,255\}$; $n$ и $m$~--- размер растра по горизонтали и 


вертикали соответственно.





Результирующая матрица цветов пикселей изображения:


$$


\mathrm{I}\_\mathrm{L}=A_1\cup A_2\cup A_3\,;


$$


изображения левого и правого кадров:


$$


\mathrm{IL} = \mathrm{I}\_\mathrm{L}_1 \ldots \mathrm{I}\_\mathrm{L}_i\,,\enskip 


\mathrm{IR} = \mathrm{I}\_\mathrm{R}_1\ldots \mathrm{I}\_\mathrm{R}_i\,,


$$


где $i$~--- этапы съемки объекта.


      


\subsection{Аппроксимация цветового пространства изображений} %2.2





В процессе работы с изображениями, получаемыми с разных ракурсов при помощи 


цифровой фотокамеры, возможно не идентичное распознавание одного и того же цвета 


пикселей, что\linebreak мешает однозначному определению ОТ. Для решения задачи\linebreak идентичного 


распознавания используется цветовая индексная\linebreak палитра. 





Цветовая индексная палитра представляет собой совокуп-\linebreak ность пикселей, индексы, 


которых соответствуют присвоенным им синхронизированным цветовым данным левого 


и правого\linebreak изоб\-ра\-жений:


$$


\mathrm{CIP}^k=\begin{bmatrix}


R_1 & G_1 & \ldots & B_1\\


R_2 & G_2 & \ldots & B_2\\


\ldots & \ldots & \ldots & \ldots\\


R_l & G_l &\ldots & B_l


\end{bmatrix}\,,


$$


где $l$~--- количество элементов цветовых индексных палитр CIP$^k$; $k$~--- индекс 


цветовой индексной палитры.





В процессе присвоения индексу цветовых данных, цвета, имеющие небольшое различие в 


цветовых характеристиках, будут игнорироваться, т.\,е.\ не различимым цветам с 


определенным шагом различия будут присвоены одинаковые индексы. Определение 


различия цветов проводится путем определения евклидового расстояния между 


цветовыми характеристиками: 


$$


\mathrm{Cl} =\sqrt{(R_1-R_2)^2+(G_1-G_2)^2+(B_1-B_2)^2}\,,


$$


где $R$~--- значение интенсивности красного цвета; $G$~--- значение интенсивности 


зеленого цвета; $B$~--- значение интенсивности синего цвета.





Таким образом, реальное количество цветов, представленное на исходном изображении, 


будет сокращено до определенного размера. Это позволит однозначно производить 


распознавание цветовых данных в результате небольшого изменения цвет\-ности пикселей


по соответствующим индексам. 





Информационные данные, используемые в процессе создания цветовой индексной 


палитры, следующие:


$R, G, B$~--- цветовые данные совокупностей пикселей изображений;


CIPL, CIPR~--- совокупность значений цветовой индексной палитры для левых и правых 


изображений.





\subsection{Сегментация изображений по цветовым характеристикам пикселей} %2.3





После аппроксимации цветового пространства необходимо\linebreak сегментировать изображения 


по каждому цвету, содержащемуся в цветовой индексной палитре. Процесс сегментации 


заключается в разграничении изображения на области пикселей с соответствующими 


цветовыми характеристиками. 





Пусть $f(x,y)$~--- функция, описывающая анализируемое изоб\-ра\-жение; $X$~--- конечное 


множество точек, на котором определена функция~$f(x,y)$; $S = \{S_1, S_2, \ldots , S_k\}$~--- 


разбиение~$X$ на $k$~об\-ла\-стей~$S_i$, где $i = 1,2, \ldots , k$; Cp~--- предикат, заданный 


на множестве~$S$ и принимающий истинные значения тогда и только тогда,\linebreak когда любая 


пара точек~$(x,y)$ из каждого подмножества~$S_i$\linebreak удовле\-тво\-ря\-ет некоторому критерию 


одно\-род\-ности этого подмножества.





Сегментацией изображения~$f(x,y)$ по предикату Cp называется разбиение $S = \{S_1,S_2, 


\ldots , S_k\}$, удовлетворяющее условиям:


\begin{enumerate}[(1)]


\item $\bigcup\limits_{i=1}^k S_i=X$;


\item $S_i \bigcap S_j=\varnothing$ для любых $i\not = j$;


\item  $\mathrm{Cp}(S_i) = \mbox{<<истина>>}$ для любого~$i$;


\item $\mathrm{Cp}(S_i  \bigcup S_j) = \mbox{<<ложно>>}$ для любых $i\not= j$.


\end{enumerate}





За счет использования, полученной на предыдущем этапе, цветовой индексной палитры 


осуществляется распознавание схожих областей на изображениях. Для определения 


иден\-тич\-ности областей на разных снимках вычисляется их положение относительно 


маркера сцены. Сравнение идентичных областей на снимках, полученных на разных 


временных этапах, позволяет проводить анализ возможных изменений формы 


по\-верх\-ности\linebreak объекта. 





\subsection{Структурирование сегментированных областей} %2.4





Следующим шагом является структурирование областей.\linebreak Это необходимо для выбора 


характеристических точек внутри об\-ласти, которые используются для вычисления 


трехмерных координат, а также для определения изменений поверхности 


объекта, заключающихся в изменении структуры прямоугольников, а значит и смещении 


характеристических точек.





Расширим математическое описание сегментации изображения~$f(x,y)$ условиями структурирования сегментированных об-\linebreak ластей:


$T = \{T_1, T_2, \ldots , T_n\}$~--- разбиение~$S_i$ на~$n$ областей~$T_j$, где $j = 1, 2, \ldots , 


n$.





Тогда структурированием сегментированной области~$S_i$ называется разбиение $T = 


\{T_1, T_2, \ldots , T_n\}$, удовлетворяющее условиям:


\begin{enumerate}[(1)]


\item $\bigcup\limits_{j=1}^k T_j=S_i$;


\item $T_i\bigcap T_j=\varnothing$  для любых $i \not= j$.


\end{enumerate}





В качестве структурирующего элемента выбран пря\-мо\-уголь\-ник, так как он прост в 


описании и не требуется сложных ал\-го\-рит\-мов для реализации вычислений, связанных с 


ним. Ко\-ли\-чество структурирующих прямоугольников (СП) выбирается относи\-тель\-но 


размеров области на основании эмпирических данных~\cite{6-7}:


$$


\mathrm{SPL}=\begin{bmatrix}


\mathrm{tl}x_1 & \mathrm{tl}y_1 & \mathrm{br}x_1 & \mathrm{br}y_1\\


\mathrm{tl}x_2 & \mathrm{tl}y_2 & \mathrm{br}x_2 & \mathrm{br}y_2\\


\ldots &\ldots &\ldots &\ldots\\


\mathrm{tl}x_n & \mathrm{tl}y_n & \mathrm{br}x_n & \mathrm{br}y_n\\


\end{bmatrix}\,,


$$


где $n$~--- количество структурирующих элементов; tl$x$, tl$y$, br$x$, br$y$~--- 


координаты вершин верхней левой (TopLeft) и нижней правой (BottomRight) 


соответственно.





Структурирующие прямоугольники близко расположенные к границам области, могут 


содержать пиксели других областей,\linebreak такие прямоугольники являются не 


информативными. Использование таких прямоугольников для дальнейшего анализа 


может привести к ошибкам вычисления трехмерных координат,\linebreak поэтому они 


исключаются из рассмотрения. Для определения информативности структурирующих 


прямоугольников должно выполняться условие равенства количества пикселей дан\-но\-го\linebreak\vspace*{-12pt}


\pagebreak





\noindent 


цве\-та сегментированной области в структурирующем прямо-\linebreak угольнике и общего числа 


пикселей в нем:


$$


\sum\limits_{i=1}^n \mathrm{Cp}(\mathrm{pixel}_i)=\sum\limits_{j=1}^n \mathrm{pixel}_j\,,


$$


где Cp~--- предикат, принимающий истинные значения тогда и только тогда, когда 


точка (пиксель) удовлетворяет некоторому критерию однородности, в данном случае цвету 


области; $n$~--- количество пикселей в СП.





Формируется матрица следующего вида:


$$


\mathrm{OS}=


\begin{bmatrix}


\mathrm{tl}x_1 & \mathrm{tl}y_1 & \mathrm{br}x_1 & \mathrm{br}y_1 & s\\


\mathrm{tl}x_2 & \mathrm{tl}y_2 & \mathrm{br}x_2 & \mathrm{br}y_2 & s\\


\ldots &\ldots &\ldots &\ldots &\ldots\\


\mathrm{tl}x_n & \mathrm{tl}y_n & \mathrm{br}x_n & \mathrm{br}y_n & s\\


\end{bmatrix}\,,


$$


где  $s$~--- информативное состояние ОП, $s\in \{0,\,1\}$. 





Структурирование проводится в несколько этапов. Первым из них является выбор области 


на исходном изображении.\linebreak\vspace*{-12pt}





\begin{floatingfigure}{57mm} %fig2-7


%\vspace*{6pt}


\begin{center}


\mbox{%


\epsfxsize=50.669mm


\epsfbox{arh-2.eps}


}


 \end{center}


 \vspace*{-9pt}


\Caption{Структурирование сегментированной области на исходном изображении


\label{f2-7}}


\vspace*{9pt}


\end{floatingfigure}





\noindent


 В~качестве исходного изоб\-ра\-же\-ния используется снимок,\linebreak


от\-но\-си\-тель\-но которого вы-\linebreak чис\-ля\-ют\-ся возможные изме-\linebreak


нения об\-ласти на сле\-ду\-ющем по 


временному показателю\linebreak сним\-ке. Итоговым изобра\-же\-ни\-ем будет следующее после 


исходного. Исходя из размеров области, проводится разбиение на структурирующие 


прямоугольники, как показа-\linebreak но на рис.~2.





      


На итоговом изображении осуществляется поиск иден-\linebreak тичной области по цветовым 


характеристикам и положе-\linebreak нию относительно маркера. Так как сцена с маркерами и 


ре\-гист\-ри\-ру\-ющие приборы статичны, изменениям подвержен только натурный объект, то 


осуществляем перенос структурирующей сетки на итоговое изображение. В~соответствии 


с рис.~3,\,\textit{а}, координаты положения центра структурирующей области на итоговом 


изоб\-ра\-же\-нии соответствуют координатам центра на исходном.


                     


\begin{figure} %fig3-7


%\vspace*{6pt}


\begin{center}


\mbox{%


\epsfxsize=107.174mm


\epsfbox{arh-3.eps}


}


 \end{center}


 \vspace*{-9pt}


\Caption{Начальное~(\textit{а}) и финальное~(\textit{б})


 структурирование сегментированной области на итоговом  изображении


\label{f3-7}}


\end{figure}


      


      Следующим этапом является проверка соответствия 


структурирующих прямоугольников на изображениях. 





Условием 


соответствия является наличие заданного количества информативных СП в 


пределах области. Исходя из этого осуществляется изменение размеров СП, 


таким образом, что бы количество информативных прямоугольников 


оставалось неизменным. На рис.~3,\,\textit{б} представлен финальный этап 


структурирования сегментированной области итогового изображения.


                        


     


В процессе структурирования сегментированных областей используются следующие 


данные:


\begin{itemize}


\item SEGL, SEGR~--- совокупность сегментированных областей на левом и правом 


изображениях соответственно;


\item $[\mathrm{SPL}]_j$, [SPR]$_j$~--- совокупность структурирующих прямоугольников для $j$-й 


области на левом и правом изображении соответственно.


\end{itemize}





\subsection{Фиксация опорных точек на плоских изображениях} %2.5





Завершающим этапом выбора ОТ на двумерных изображениях является их фиксация.





Метод фиксации опорных точек на изображениях заключается в выборе определенных 


пикселей или областей пикселей на начальных изображениях и поиске аналогичных на 


следующих снимках. Поиск и сравнение координат ОТ позволяют определять изменения 


формы поверхности натурного объекта. Таким образом, для изменения трехмерной 


модели объекта необходимы только новые значения координат точек. Это позволяет 


снизить объем передаваемой информации и повысить скорость моделирования объекта.





В качестве опорных точек будут использоваться пиксели, расположенные на пересечении 


диагоналей структурирующего элемента:


$$


\mathrm{pixel} \left[\left(\fr{x_2+x_1}{2}\right),\, \left(\fr{y_2+y_1}{2}\right)\right]\,.


$$





Далее проводится проверка равенства координат ОТ. Координаты, изменившие свое 


положение относительно данных на предыдущем изображении, сохраняются и 


передаются для вычислений трехмерных координат и изменений модели натурного 


объекта. Для передачи используются следующие данные:


\begin{itemize}


\item SP~--- исходные координаты точек, изменивших свое положение;


\item


FP~--- итоговые координаты точек, изменивших свое положение.


\end{itemize}





\subsection{Вычисление координат опорных точек в трехмерном 


пространстве по двумерным координатам} %2.6





Вычисление координат опорных точек в трехмерном про-\linebreak стран\-ст\-ве по двумерным 


координатам на плоских изображениях производится путем аналитической обработки 


исходных данных в соответствии с математической моделью, основанной на принципах 


фотограмметрии, проекционной геометрии и стереоскопии~\cite{2-7, 7-7}. В~связи с этим 


фото- или видеосъемка производится с двух разных ракурсов (камер), разнесенных на 


величину стереобазиса. В~качестве исходных данных используются двумерные 


координаты опорных точек, идентифицируемых на плоских изоб\-ра\-же\-ниях.





\subsection{Построение и преобразование модели} %2.7


Трехмерная модель объекта представлена воксельной мо-\linebreak делью. Данный выбор 


обусловлен рядом преимуществ:


\begin{itemize}


\item простота выполнения булевых операций пересечения и объединения объектов;


\item возможность естественной трансформации между несколькими объектами 


(морфинга);


\item снятие ограничений на геометрическую сложность сцены. Единственным и 


универсальным фактором, влияющим на сложность моделируемых объектов, является 


воксельное разрешение.


\end{itemize}





Воксельная модель представляется совокупностью вокселей, которые являются кубами в 


трехмерном пространстве и могут быть описаны следующим образом:


$$


V_i=\{x,y,z,r,g,b,t\}\,,\quad i=1,2,\ldots , n\,,


$$





\begin{floatingfigure}{55mm} %fig5-7


%\vspace*{6pt}


\begin{center}


\mbox{%


\epsfxsize=50.434mm


\epsfbox{arh-5.eps}


}


 \end{center}


 \vspace*{-9pt}


\Caption{Декомпозиция области морфинга на плоскости


\label{f5-7}}


\vspace*{9pt}


\end{floatingfigure}





\noindent


где $x,y,z$~--- координаты $i$-го вокселя в трехмерном про-\linebreak странстве; $r,g,b$~--- цветовые 


координаты $i$-го вокселя в пространстве RGB; $t$~--- степень прозрачности $i$-го 


вокселя, $t=$\linebreak $=0,\ldots , 1$; $n$~--- количество\linebreak вокселей в модели.





Для построения модели,\linebreak пространство между опорны-\linebreak ми точками интерполируется\linebreak 


вокселями при помощи метода воксельной аппроксимации\linebreak плоскостями~\cite{8-7, 9-7}. 


Для этого поверхность модели, ограни-\linebreak ченная ОТ разбивается на об-\linebreak ласти, для каждой ОТ 


область ограничивается соседними опорными точками. В~соответствии с рис.~4, для 


опорной точки~$O$ об-\linebreak ласть ограничивается плоскостями $(O_1, O, O_2), (O_2, O, O_3), 


\ldots$\linebreak $\ldots , (O_7, O, O_8)$, проходящими через нее и соседние с ней ОТ $(O_1,$\linebreak $ O_2, O_3, O_4, 


O_5, O_6, O_7, O_8)$. 








       


Поверхность можно представить в виде совокупности плоскостей:


\pagebreak





\noindent


$$


S =(O_1 O_2 O_3)\cup (O_3O_4O_5)\cup\ldots\cup(O_{n-2}O_{n-1}O_n)\,,


$$


где $O_i$~--- опорные точки модели.





Таким образом, для построения воксельной модели для каж\-дой ОТ необходимо:


\begin{enumerate}[1.]


\item  Определить соседние опорные точки.





Выбор соседних ОТ будет осуществляться с учетом евклидова расстояния между ними на 


плоскости:


%\noindent


%\setcounter{equation}{0}


\begin{equation}


d=\sqrt{(x_1-x_2)^2+(y_1-y_2)^2}\,,


\label{e11}


\end{equation}


где $x_1$, $y_1$~--- координаты ОТ, изменившей свое положе-\linebreak ние~$XOY$; $x_2$, $y_2$~--- 


координаты рассматриваемой ОТ на плоскости~$XOY$.





Использование только координат на плоскости достаточно, учитывая тот факт, что 


рассматриваются выпуклые поверхности. В~случае если имеется поверхность другого 


вида, она должна аппроксимироваться до выпуклой. Используя формулу~(\ref{e11}), 


находятся все расстояния между ОТ и со\-став\-ля\-ет\-ся матрица расстояний вида:


$$


S=\begin{bmatrix}


0 & s_{12} &\ldots & s_{1n}\\


s_{21} & 0 &\ldots & s_{2n}\\


\ldots &\ldots &\ldots &\ldots\\


s_{n1} & s_{n2} &\ldots & 0\\


\end{bmatrix}\,,


$$


где $n$~--- количество ОТ.


При этом $s_{ij}=s_{ji}$.


Затем необходимо отсортировать ОТ по возрастанию расстояний друг для друга. 


В~результате получим матрицу размерностью $n \times  (n-1)$, в строках которой 


содержатся индексы ОТ, отсортированные по удалению от ОТ с индексом, равным 


индексу столбца.





\item  Вычислить уравнения каждой из плоскостей~\cite{8-7}.





Для этого необходимо вычислить коэффициенты общего\linebreak уравнения плоскости, 


проходящей через три точки в пространстве:


$$


Ax+By+Cz+D=0\,,


$$


которые равны соответственно:


\begin{align*}


A &=y_1(z_2-z_3)+y_2(z_3-z_1)+y_3(z_1-z_2)\,;\\


B&= x_1(z_2-z_3)+x_2(z_3-z_1)+x_3(z_2-z_2)\,;\\


C &= x_1(y_2-y_3)+x_2(y_3-y_1)+x_3(y_1-y_2)\,;\\


D &= x_1(y_2z_3-y_3z_2)+x_2(y_3z_1-y_1z_3)+x_3(y_1z_2-y_2z_1)\,.


\end{align*}


\item По уравнениям вычислить положения промежуточных точек.


\item Проверить, попадают ли точки в область, ограниченную треугольниками.


\end{enumerate}





Подробное описание алгоритма приведено в~\cite{9-7}.





В результате получим воксельную модель объекта в трехмерном пространстве


$$


V=\sum\limits_{i=1}^n V_i\,,


$$


где $n$~--- число вокселей в модели.





Преобразование модели основывается на обработке изменений положения опорных точек. 


При изменении их положения необходимо вычислить смещение окружающих данную 


опорную точку вокселей в определенной области. После этого для каж\-дой ОТ, 


изменившей свое положение осуществить воксельную аппроксимацию по описанному 


выше алгоритму.





\section{Заключение}





Таким образом, был разработан метод регистрации мор\-фин\-га объектов в трехмерном 


пространстве. Рассмотренный метод не заключает в себе недостатков существующих: 


способ регистрации исходных данных является бесконтактным и не требует 


специализированных приборов, сравнительно дешев и не требует высокой квалификации 


пользователей.





Кроме того, данный метод позволяет ускорить процесс создания трехмерных моделей 


реальных объектов, а также позволяет осуществлять сбор, хранение и восстановление 


информации об изменениях поверхности объекта.





{\frenchspacing\small


%\baselineskip=10.5pt


\begin{thebibliography}{9}





\bibitem{1-7}


\Au{Дамьяновски В.}


CCTV. Библия охранного телевидения.~--- 2-е изд.~--- М.: Ай-Эс-Эс Пресс, 2003. 





\bibitem{2-7}


\Au{Архипов О.\,П. Маньяков Ю.\,А., Сиротинин~Д.\,О.}


Вычисление пространственных координат опорных точек морфинга по плоским 


изоб\-ра\-же\-ни\-ям~// Изв.\ ОрелГТУ, 2008.  №\,3/271(546). С.~18--24.





\bibitem{3-7}


\Au{Валюс Н.\,А.}


Стереоскопия.~--- М.: Изд-во Акад. наук СССР, 1962.





\bibitem{4-7}


\Au{Архипов О.\,П., Маньяков Ю.\,А., Сиротинин~Д.\,О.}


Технология морфинга трехмерных моделей на основе данных натурного 


эксперимента~// Изв.\ ОрелГТУ, 2009. №\,1. С.~9--13.





\bibitem{5-7}


\Au{Маньяков Ю.\,А., Сиротинин Д.\,О.}


Автоматизация процесса моделирования поверхности 3D объекта в ограниченном 


трехмерном пространстве~// Компьютерные науки и технологии. Ч.~2.~--- Белгород: 


ГиК, 2009.~--- С.~194--198.





\bibitem{6-7}


\Au{Архипов О.\,П., Маньяков Ю.\,А., Сиротинин~Д.\,О.}


Вариант построения сетки опорных точек по цветовым данным растрового 


изображения~// Изв.\ ОрелГТУ, %Информационные системы и технологии, 


2008.  №\,4/272(550). С.~34--39.





\bibitem{7-7}


\Au{Архипов О.\,П., Маньяков Ю.\,А., Сиротинин~Д.\,О.}


Разработка и исследование информационной технологии морфинга 3D объекта на 


основе данных натурного эксперимента: Отчет о НИР (промежуточный)~/ ОФ ИПИ РАН.~--- Орел, 2008. 129~с.





\bibitem{8-7}


\Au{Архипов О.\,П., Маньяков Ю.\,А., Сиротинин Д.\,О.}


Разработка и исследование информационной технологии морфинга 3D объекта на 


основе данных натурного эксперимента: Отчет о НИР (промежуточный)~/ ОФ ИПИ РАН.~--- Орел, 2009.  97~с. 





\bibitem{9-7}


\Au{Маньяков Ю.\,А.}


Технология регистрации поведения объектов в трехмерном пространстве~// 


Информационные технологии в науке, образовании и производстве. ИТНОП-2010.~--- 


Орел: ОрелГТУ, 2010. Т.~3. С.~182--186.


\end{thebibliography}


}


